\documentclass[12pt,a4paper]{article}

% Pacotes essenciais
\usepackage[utf8]{inputenc}
\usepackage[T1]{fontenc}
\usepackage[english]{babel}
\usepackage{amsmath, amssymb, amsthm}
\usepackage{geometry}
\usepackage{graphicx}
\usepackage{xcolor}
\usepackage{listings}
\usepackage{hyperref}
\usepackage{booktabs}
\usepackage{array}
\usepackage{fancyhdr}
\usepackage{titlesec}
\usepackage{tcolorbox}
\usepackage{enumitem}
\usepackage{float}

% Geometria da página
\geometry{
  top=2.5cm, bottom=2.5cm,
  left=3cm, right=2.5cm
}

% Cores do projeto
\definecolor{myblue}{RGB}{30, 90, 160}
\definecolor{mygreen}{RGB}{0, 128, 80}
\definecolor{mygray}{RGB}{245, 245, 245}
\definecolor{myorange}{RGB}{210, 100, 10}
\definecolor{myred}{RGB}{180, 30, 30}
\definecolor{codegray}{RGB}{240, 240, 240}
\definecolor{codecomment}{RGB}{100, 130, 100}
\definecolor{codestring}{RGB}{160, 40, 40}
\definecolor{codekeyword}{RGB}{30, 80, 160}

% Cabeçalho e rodapé
\pagestyle{fancy}
\fancyhf{}
\fancyhead[L]{\textcolor{myblue}{\small \textbf{My Chance} — Documentação Técnica}}
\fancyhead[R]{\textcolor{gray}{\small Algoritmos de Matching}}
\fancyfoot[C]{\textcolor{gray}{\thepage}}
\renewcommand{\headrulewidth}{0.4pt}

% Estilo de títulos
\titleformat{\section}
  {\Large\bfseries\color{myblue}}
  {\thesection.}{0.8em}{}[\color{myblue}\titlerule]

\titleformat{\subsection}
  {\large\bfseries\color{myblue!80!black}}
  {\thesubsection.}{0.6em}{}

\titleformat{\subsubsection}
  {\normalsize\bfseries\color{myblue!60!black}}
  {\thesubsubsection.}{0.5em}{}

% Caixas coloridas
\tcbuselibrary{skins, breakable}

\newtcolorbox{infobox}[1][]{
  enhanced, breakable,
  colback=myblue!5!white,
  colframe=myblue,
  fonttitle=\bfseries,
  title={#1},
  boxrule=0.8pt,
  left=6pt, right=6pt, top=4pt, bottom=4pt
}

\newtcolorbox{formulabox}[1][]{
  enhanced,
  colback=mygreen!5!white,
  colframe=mygreen,
  fonttitle=\bfseries,
  title={#1},
  boxrule=0.8pt,
  left=6pt, right=6pt, top=4pt, bottom=4pt
}

\newtcolorbox{examplebox}[1][]{
  enhanced, breakable,
  colback=myorange!5!white,
  colframe=myorange,
  fonttitle=\bfseries,
  title={#1},
  boxrule=0.8pt,
  left=6pt, right=6pt, top=4pt, bottom=4pt
}

\newtcolorbox{warnbox}[1][]{
  enhanced,
  colback=myred!5!white,
  colframe=myred,
  fonttitle=\bfseries,
  title={#1},
  boxrule=0.8pt,
  left=6pt, right=6pt, top=4pt, bottom=4pt
}

% Estilo de código
\lstdefinestyle{pythonstyle}{
  backgroundcolor=\color{codegray},
  commentstyle=\color{codecomment}\itshape,
  keywordstyle=\color{codekeyword}\bfseries,
  stringstyle=\color{codestring},
  basicstyle=\ttfamily\small,
  breakatwhitespace=false,
  breaklines=true,
  captionpos=b,
  keepspaces=true,
  numbers=left,
  numbersep=6pt,
  numberstyle=\tiny\color{gray},
  showspaces=false,
  showstringspaces=false,
  showtabs=false,
  tabsize=4,
  frame=single,
  framerule=0.5pt,
  rulecolor=\color{gray!40},
  language=Python
}

\lstset{style=pythonstyle}

% Hiperlinks
\hypersetup{
  colorlinks=true,
  linkcolor=myblue,
  urlcolor=myblue,
  citecolor=myblue
}

% --------------------------------------------------
\begin{document}

% Capa
\begin{titlepage}
  \centering
  \vspace*{2cm}
  {\Huge\bfseries\color{myblue} My Chance\par}
  \vspace{0.3cm}
  {\large\itshape Onde seu perfil profissional é o mais importante\par}
  \vspace{1.5cm}
  \rule{\linewidth}{1.5pt}
  \vspace{0.3cm}
  {\LARGE\bfseries Documentação Técnica\par}
  \vspace{0.2cm}
  {\Large\color{myblue!70!black} Módulo de Algoritmos de Matching\par}
  \vspace{0.2cm}
  \rule{\linewidth}{1.5pt}
  \vspace{1.2cm}
  {\large \textbf{Item 4 — Arquitetura e Principais Tecnologias}\par}
  \vspace{1cm}
  \begin{tcolorbox}[colback=myblue!8, colframe=myblue, width=0.85\linewidth]
    \centering
    \textbf{Algoritmos documentados neste arquivo:}\\[6pt]
    1. Pontuação Ponderada por Competências\\
    2. Similaridade por Cosseno (Cosine Similarity)
  \end{tcolorbox}
  \vspace{1.5cm}
  {\large Equipe My Chance — Disciplina de MVP\par}
  \vspace{0.4cm}
  {\normalsize Davi Lima de Medeiros \quad Israel Buriti Galvão\par}
  {\normalsize João Victor Cosme Melo \quad José do Bomfim Truta Neto\par}
  {\normalsize Lucas Leonardo Mayer Almeida \quad Lucas Queiroz Correia\par}
  \vfill
  {\normalsize Campina Grande, Paraíba — 2026\par}
\end{titlepage}

% Sumário
\tableofcontents
\newpage

% --------------------------------------------------
\section{Introdução e Contexto}

O \textbf{My Chance} é uma plataforma de recrutamento às cegas e invertido, cujo objetivo central é
reduzir a influência de vieses inconscientes na triagem de candidatos. Em vez de candidatos
se candidatarem a vagas, \emph{recrutadores} buscam ativamente perfis no banco de dados,
enxergando apenas competências técnicas e experiências — os dados pessoais são revelados
somente após o aceite mútuo de uma entrevista.

Para que esse mecanismo funcione com qualidade e justiça, o sistema precisa de um
\textbf{motor de matching inteligente}: um conjunto de algoritmos que ordene candidatos por
compatibilidade com uma vaga de forma precisa, transparente e escalável.

Este documento descreve, de forma didática e detalhada, os dois algoritmos propostos para
compor esse motor, incluindo fundamentos matemáticos, implementação em Python e
orientações de integração com o restante da arquitetura (Streamlit, SQLite/PostgreSQL,
Docker).

\begin{infobox}[Arquitetura em Camadas]
A sugestão estratégica para o MVP é combinar os dois algoritmos em sequência:
\begin{enumerate}
  \item \textbf{Pontuação Ponderada} — filtragem rápida e eliminação de candidatos com lacunas críticas.
  \item \textbf{Cosine Similarity} — ranking principal de compatibilidade.
\end{enumerate}
\end{infobox}

\newpage

% --------------------------------------------------
\section{Algoritmo 1 — Pontuação Ponderada por Competências}

\subsection{Fundamentos Teóricos}

A \textbf{Pontuação Ponderada por Competências} (\textit{Weighted Skill Scoring}) é um método de ranqueamento baseado em somas ponderadas. É um algoritmo mais simples, mas também o mais transparente e explicável, características essenciais para um MVP que precisará convencer recrutadores.

\begin{infobox}[Intuição]
Pense como um recrutador humano preenchendo uma planilha: para cada competência
exigida pela vaga, ele atribui uma importância (peso) e avalia o nível do candidato.
A nota final é a soma dessas avaliações. Simples, auditável, justo.
\end{infobox}

\subsection{Fórmula Matemática}

Dado um candidato $C$ e uma vaga $V$ com $n$ competências, o score é:

\begin{formulabox}[Fórmula da Pontuação Ponderada]
\begin{equation}
  \text{Score}(C, V) = \sum_{i=1}^{n} w_i \cdot l_i
  \label{eq:wss}
\end{equation}
onde:
\begin{itemize}
  \item $w_i \in \mathbb{R}^+$ — peso da competência $i$ definido pela vaga;
  \item $l_i \in [0, L_{\max}]$ — nível declarado do candidato na competência $i$;
  \item $L_{\max}$ — nível máximo possível (e.g., 5 em uma escala de 1–5).
\end{itemize}
\end{formulabox}

\vspace{0.4cm}

Para normalizar o score entre 0 e 1, podemos dividir pelo score máximo teórico:

\begin{formulabox}[Score Normalizado]
\begin{equation}
  \text{Score}_{\text{norm}}(C, V) = \frac{\sum_{i=1}^{n} w_i \cdot l_i}{\sum_{i=1}^{n} w_i \cdot L_{\max}}
  \label{eq:wss_norm}
\end{equation}
\end{formulabox}

\subsection{Extensões para o My Chance}

Além da fórmula base, o sistema pode incorporar bônus e penalizações:

\begin{formulabox}[Fórmula Estendida]
\begin{equation}
  \text{Score}_{\text{ext}}(C, V) = \underbrace{\sum_{i=1}^{n} w_i \cdot l_i}_{\text{base}}
  + \underbrace{\beta_p \cdot N_p}_{\text{bônus projetos}}
  + \underbrace{\beta_r \cdot R_r}_{\text{bônus recência}}
  - \underbrace{\gamma_m \cdot M_o}_{\text{penalização obrig.}}
  \label{eq:wss_ext}
\end{equation}
onde:
\begin{itemize}
  \item $N_p$ — número de projetos reais relacionados; $\beta_p$ — bônus unitário por projeto;
  \item $R_r$ — fator de recência da experiência ($\in [0,1]$, decai com o tempo);
  \item $M_o$ — número de competências obrigatórias ausentes; $\gamma_m$ — penalização por ausência.
\end{itemize}
\end{formulabox}

\subsection{Exemplo Numérico Detalhado}

\begin{examplebox}[Exemplo — Vaga de Analista de Dados]
\textbf{Requisitos da vaga:}
\begin{center}
\begin{tabular}{lcc}
  \toprule
  Competência & Peso $(w_i)$ & Obrigatória? \\
  \midrule
  Python     & 5 & Sim \\
  SQL        & 4 & Sim \\
  Docker     & 3 & Não \\
  Power BI   & 2 & Não \\
  \bottomrule
\end{tabular}
\end{center}
\textbf{Perfil do candidato A} (nível em escala 0–5):
Python = 4, SQL = 5, Docker = 2, Power BI = 0, projetos = 2

\textbf{Cálculo:}
\begin{align*}
  \text{Score}_{\text{base}} &= (5 \cdot 4) + (4 \cdot 5) + (3 \cdot 2) + (2 \cdot 0) = 46 \\
  \text{Score}_{\text{máx}}  &= (5 + 4 + 3 + 2) \cdot 5 = 70 \\
  \text{Score}_{\text{norm}} &= \frac{46}{70} \approx 0{,}657 \\
  \text{Bônus projetos}      &= 2 \cdot 2{,}0 = +4{,}0 \\
  \text{Score}_{\text{final}} &\approx 0{,}657 + 0{,}057 \approx 0{,}714
\end{align*}
\end{examplebox}

\subsection{Implementação em Python}

\begin{lstlisting}[caption={Pontuação ponderada com bônus e penalizações},
                   label={lst:wss}]
from dataclasses import dataclass, field
from typing import Optional


@dataclass
class Vaga:
    id: int
    titulo: str
    # Ex: {"python": {"peso": 5, "obrigatoria": True, "nivel_min": 3}}
    competencias: dict


@dataclass
class Candidato:
    id: int
    nome_anonimo: str          # ID anonimo no My Chance
    # Ex: {"python": 4, "sql": 5, "docker": 2}
    competencias: dict
    projetos: int = 0          # Projetos reais no portfolio
    anos_experiencia: float = 0.0


def calcular_fator_recencia(anos: float) -> float:
    """
    Decaimento exponencial: experiencias recentes valem mais.
    Anos 0-1 -> fator ~1.0
    Anos 5   -> fator ~0.6
    Anos 10+ -> fator ~0.37
    """
    import math
    return math.exp(-0.1 * anos)


def calcular_score_ponderado(
    candidato: Candidato,
    vaga: Vaga,
    nivel_maximo: int = 5,
    bonus_por_projeto: float = 2.0,
    penalizacao_obrigatoria: float = 15.0,
) -> dict:
    """
    Calcula o score ponderado de um candidato para uma vaga.

    Retorna um dict com score bruto, normalizado e detalhes.
    """
    score_base = 0.0
    score_maximo = 0.0
    ausencias_obrigatorias = 0
    detalhes = {}

    for competencia, config in vaga.competencias.items():
        peso = config["peso"]
        obrigatoria = config.get("obrigatoria", False)
        nivel_candidato = candidato.competencias.get(competencia, 0)

        contribuicao = peso * nivel_candidato
        score_base += contribuicao
        score_maximo += peso * nivel_maximo

        detalhes[competencia] = {
            "peso": peso,
            "nivel_candidato": nivel_candidato,
            "contribuicao": contribuicao,
        }

        # Verifica ausencia de competencia obrigatoria
        nivel_minimo = config.get("nivel_min", 1)
        if obrigatoria and nivel_candidato < nivel_minimo:
            ausencias_obrigatorias += 1

    # Score normalizado [0, 1]
    score_normalizado = score_base / score_maximo if score_maximo > 0 else 0.0

    # Bonus por projetos
    bonus_projetos = (candidato.projetos * bonus_por_projeto) / score_maximo

    # Fator de recencia
    fator_recencia = calcular_fator_recencia(candidato.anos_experiencia)

    # Penalizacao por competencias obrigatorias ausentes
    penalizacao = (ausencias_obrigatorias * penalizacao_obrigatoria) / score_maximo

    # Score final
    score_final = max(
        0.0,
        (score_normalizado + bonus_projetos) * fator_recencia - penalizacao
    )

    return {
        "candidato_id": candidato.id,
        "score_base_normalizado": round(score_normalizado, 4),
        "score_final": round(score_final, 4),
        "ausencias_obrigatorias": ausencias_obrigatorias,
        "detalhes_por_competencia": detalhes,
    }


def rankear_candidatos(candidatos: list[Candidato],
                        vaga: Vaga) -> list[dict]:
    """Ordena candidatos por score decrescente."""
    resultados = [
        calcular_score_ponderado(c, vaga) for c in candidatos
    ]
    return sorted(resultados, key=lambda x: x["score_final"],
                  reverse=True)


# ------------------------------------------------------------------
# Exemplo de uso
# ------------------------------------------------------------------
if __name__ == "__main__":
    vaga_analista = Vaga(
        id=1,
        titulo="Analista de Dados",
        competencias={
            "python": {"peso": 5, "obrigatoria": True,  "nivel_min": 3},
            "sql":    {"peso": 4, "obrigatoria": True,  "nivel_min": 2},
            "docker": {"peso": 3, "obrigatoria": False, "nivel_min": 1},
            "powerbi":{"peso": 2, "obrigatoria": False, "nivel_min": 1},
        }
    )

    candidatos = [
        Candidato(1, "CAND-001",
                  {"python": 4, "sql": 5, "docker": 2, "powerbi": 0},
                  projetos=2, anos_experiencia=1.0),
        Candidato(2, "CAND-002",
                  {"python": 5, "sql": 4, "docker": 4, "powerbi": 3},
                  projetos=5, anos_experiencia=3.0),
        Candidato(3, "CAND-003",
                  {"python": 2, "sql": 1, "docker": 0, "powerbi": 5},
                  projetos=0, anos_experiencia=0.5),
    ]

    ranking = rankear_candidatos(candidatos, vaga_analista)
    for pos, r in enumerate(ranking, 1):
        print(f"{pos}. {r['candidato_id']} "
              f"| Score: {r['score_final']:.3f} "
              f"| Ausencias obrig.: {r['ausencias_obrigatorias']}")
\end{lstlisting}

\subsection{Vantagens e Limitações}

\begin{table}[H]
\centering
\caption{Análise da Pontuação Ponderada no projeto}
\begin{tabular}{|p{7cm}|| p{7cm}|}
  \toprule
  \textbf{Vantagens} & \textbf{Limitações} \\
  \hline
  Extremamente simples de implementar & Pesos iniciais são subjetivos (recrutador define) \\
  \hline
  Completamente transparente e auditável & Não captura similaridade estrutural entre perfis \\
  \hline
  Fácil de explicar em apresentação & Escalas declaradas pelo candidato podem ser inconsistentes \\
  \hline
  Permite penalização por gaps críticos & Não detecta combinações sinérgicas de competências \\
  \hline
  Funciona mesmo sem histórico de dados & Ranking linear pode ser simplista para vagas complexas \\
  \bottomrule
\end{tabular}
\end{table}

\newpage

% --------------------------------------------------
\section{Algoritmo 2 — Matching por Similaridade Vetorial (Cosine Similarity)}

\subsection{Fundamentos Teóricos}

A \textbf{Similaridade por Cosseno} (\textit{Cosine Similarity}) é uma medida de similaridade
entre dois vetores no espaço euclidiano. É amplamente utilizada em sistemas de
recomendação, recuperação de informação e processamento de linguagem natural por
uma razão elegante: ela captura a \emph{direção} dos vetores, não apenas suas magnitudes.

\begin{infobox}[Intuição Geométrica]
Imagine cada competência como uma dimensão em um espaço vetorial. Tanto o candidato
quanto a vaga são representados como pontos nesse espaço. A \textit{Cosine Similarity} mede o
\textbf{ângulo} entre esses dois pontos — quanto menor o ângulo, mais similares os perfis.
Isso significa que um candidato com todas as competências no nível 4 pode ser tão
compatível com uma vaga quanto um com todas no nível 2, se a \emph{proporção} for similar.
\end{infobox}

\subsection{Fórmula Matemática}

Dados dois vetores $\mathbf{A}$ (vaga) e $\mathbf{B}$ (candidato) no $\mathbb{R}^n$:

\begin{formulabox}[Similaridade por Cosseno]
\begin{equation}
  \cos(\theta) = \frac{\mathbf{A} \cdot \mathbf{B}}{\|\mathbf{A}\| \cdot \|\mathbf{B}\|}
  = \frac{\sum_{i=1}^{n} A_i B_i}{\sqrt{\sum_{i=1}^{n} A_i^2} \cdot \sqrt{\sum_{i=1}^{n} B_i^2}}
  \label{eq:cosine}
\end{equation}
onde $\theta$ é o ângulo entre os vetores e o resultado pertence ao intervalo $[-1, 1]$.
\end{formulabox}

\vspace{0.4cm}

Para o caso de competências (valores não negativos), o resultado estará sempre em $[0, 1]$:

\begin{table}[H]
\centering
\caption{Interpretação da Cosine Similarity}
\begin{tabular}{c|c|c}
  \toprule
  $\cos(\theta)$ & $\theta$ & Interpretação \\
  \midrule
  $1{,}00$            & $0°$   & Perfis idênticos \\
  $0{,}85 - 0{,}99$  & $< 30°$ & Alta compatibilidade \\
  $0{,}60 - 0{,}84$  & $30°-53°$ & Compatibilidade moderada \\
  $0{,}30 - 0{,}59$  & $> 53°$ & Baixa compatibilidade \\
  $0{,}00$            & $90°$  & Sem competências em comum \\
  \bottomrule
\end{tabular}
\end{table}

\subsection{Construção dos Vetores no My Chance}

O passo mais importante é a \textbf{vetorização}: converter perfis textuais em vetores numéricos.

\subsubsection{Vocabulário de Competências}

Define-se um vocabulário global $\mathcal{V} = \{c_1, c_2, \ldots, c_m\}$ com todas as
competências cadastradas no sistema. Cada perfil vira um vetor $\mathbf{v} \in \mathbb{R}^m$,
onde $v_i$ representa o nível naquela competência (ou 0 se ausente).

\subsubsection{Exemplo de Vetorização}

Considere $\mathcal{V} = \{\text{Python}, \text{SQL}, \text{Docker}, \text{React}\}$:

\begin{examplebox}[Vetorização e Cálculo de Similaridade]
\textbf{Vaga:} Python=1, SQL=1, Docker=1, React=0 $\Rightarrow \mathbf{V} = (1, 1, 1, 0)$\\
\textbf{Candidato A:} Python=1, SQL=1, Docker=0, React=1 $\Rightarrow \mathbf{C_A} = (1, 1, 0, 1)$\\
\textbf{Candidato B:} Python=5, SQL=4, Docker=4, React=0 $\Rightarrow \mathbf{C_B} = (5, 4, 4, 0)$\\[8pt]

\textbf{Cálculo para Candidato A:}
\[
  \cos(\theta_A) = \frac{(1)(1)+(1)(1)+(1)(0)+(0)(1)}{\sqrt{3}\cdot\sqrt{3}} = \frac{2}{3} \approx 0{,}667
\]

\textbf{Cálculo para Candidato B:}
\[
  \cos(\theta_B) = \frac{(1)(5)+(1)(4)+(1)(4)+(0)(0)}{\sqrt{3}\cdot\sqrt{57}} = \frac{13}{\sqrt{171}} \approx 0{,}994
\]

Candidato B é mais compatível, mesmo usando escala 0–5 vs 0–1.
\end{examplebox}

\subsection{Versão Ponderada da Cosine Similarity}

É possível incorporar os pesos da vaga diretamente na similaridade:

\begin{formulabox}[Cosine Similarity Ponderada]
\begin{equation}
  \cos_w(\theta) = \frac{\sum_{i=1}^{n} w_i \cdot A_i \cdot B_i}
  {\sqrt{\sum_{i=1}^{n} w_i \cdot A_i^2} \cdot \sqrt{\sum_{i=1}^{n} w_i \cdot B_i^2}}
  \label{eq:cosine_w}
\end{equation}
\end{formulabox}

Isso combina a expressividade geométrica da Cosine Similarity com os pesos de importância
definidos pela vaga.

\subsection{Implementação em Python}

\begin{lstlisting}[caption={Cosine Similarity para matching de vagas e candidatos},
                   label={lst:cosine}]
import numpy as np
import pandas as pd
from typing import Optional


def vetorizar_perfil(competencias_perfil: dict,
                     vocabulario: list[str]) -> np.ndarray:
    """
    Converte um dicionario de competencias em vetor numpy.

    Parametros
    ----------
    competencias_perfil : {"python": 4, "sql": 3, ...}
    vocabulario         : lista ordenada de todas as competencias

    Retorna
    -------
    np.ndarray de tamanho len(vocabulario)
    """
    return np.array([
        float(competencias_perfil.get(c, 0)) for c in vocabulario
    ])


def cosine_similarity(a: np.ndarray, b: np.ndarray,
                      pesos: Optional[np.ndarray] = None) -> float:
    """
    Calcula a similaridade por cosseno entre dois vetores.
    Suporta versao ponderada se 'pesos' for fornecido.
    """
    if pesos is not None:
        # Versao ponderada
        num = np.sum(pesos * a * b)
        den = (np.sqrt(np.sum(pesos * a**2)) *
               np.sqrt(np.sum(pesos * b**2)))
    else:
        num = np.dot(a, b)
        den = np.linalg.norm(a) * np.linalg.norm(b)

    if den == 0:
        return 0.0
    return float(num / den)


def construir_matriz_scores(candidatos: list[dict],
                             vaga: dict,
                             vocabulario: list[str]) -> pd.DataFrame:
    """
    Monta o ranking completo de candidatos para uma vaga.

    Parametros
    ----------
    candidatos  : lista de dicts com 'id', 'competencias'
    vaga        : dict com 'id', 'competencias', 'pesos'
    vocabulario : lista de todas as competencias do sistema

    Retorna
    -------
    DataFrame ordenado por similaridade decrescente
    """
    vetor_vaga = vetorizar_perfil(
        vaga["competencias"], vocabulario
    )
    vetor_pesos = vetorizar_perfil(
        vaga.get("pesos", {}), vocabulario
    )
    # Se nao houver pesos, usa vetor uniforme
    if vetor_pesos.sum() == 0:
        vetor_pesos = None

    resultados = []
    for c in candidatos:
        vetor_candidato = vetorizar_perfil(
            c["competencias"], vocabulario
        )
        sim = cosine_similarity(
            vetor_vaga, vetor_candidato, vetor_pesos
        )
        resultados.append({
            "candidato_id": c["id"],
            "similaridade": round(sim, 4),
            "angulo_graus": round(
                np.degrees(np.arccos(np.clip(sim, -1, 1))), 1
            ),
        })

    df = pd.DataFrame(resultados)
    return df.sort_values("similaridade",
                          ascending=False).reset_index(drop=True)


# ------------------------------------------------------------------
# Exemplo de uso integrado
# ------------------------------------------------------------------
if __name__ == "__main__":
    vocabulario = ["python", "sql", "docker", "react",
                   "git", "powerbi", "java"]

    vaga = {
        "id": 1,
        "competencias": {"python": 1, "sql": 1, "docker": 1,
                         "git": 1},
        "pesos":        {"python": 5, "sql": 4, "docker": 3,
                         "git": 2},
    }

    candidatos = [
        {"id": "CAND-001",
         "competencias": {"python": 4, "sql": 3, "docker": 2,
                          "git": 5, "react": 1}},
        {"id": "CAND-002",
         "competencias": {"python": 5, "sql": 5, "docker": 4,
                          "git": 4}},
        {"id": "CAND-003",
         "competencias": {"react": 5, "java": 4, "powerbi": 3}},
        {"id": "CAND-004",
         "competencias": {"python": 2, "sql": 2, "docker": 1,
                          "git": 3, "powerbi": 2}},
    ]

    ranking = construir_matriz_scores(candidatos, vaga, vocabulario)
    print(ranking.to_string(index=False))
\end{lstlisting}

\subsection{Vantagens e Limitações}

\begin{table}[H]
\centering
\caption{Análise da Cosine Similarity no projeto}
\begin{tabular}{p{7cm} p{7cm}}
  \toprule
  \textbf{Vantagens} & \textbf{Limitações} \\
  \midrule
  Invariante à escala dos níveis & Não penaliza explicitamente competências obrigatórias ausentes \\
  Captura estrutura dos perfis com precisão & Sensível à esparsidade (perfis com poucas competências) \\
  Escala facilmente para milhares de candidatos & Exige vocabulário de competências bem definido \\
  Compatible com embeddings de texto no futuro & Não considera experiência temporal ou contexto \\
  Amplamente validado em sistemas de recomendação & Candidatos com listas curtas podem ser subavaliados \\
  \bottomrule
\end{tabular}
\end{table}

\newpage

% --------------------------------------------------
\section{Arquitetura Integrada dos Algoritmos}

\subsection{Visão Geral do Pipeline}

Os dois algoritmos são complementares e devem ser usados em sequência, formando
um \textbf{pipeline de matching em duas etapas}:

\vspace{0.5cm}
\begin{center}
\begin{tabular}{|c|c|c|c|}
  \hline
  \textbf{Etapa} & \textbf{Algoritmo} & \textbf{Função} & \textbf{Output} \\
  \hline
  1 & Pontuação Ponderada & Filtragem rápida & Lista reduzida \\
  \hline
  2 & Cosine Similarity & Ranking principal & Score de 0 a 1 \\
  \hline
\end{tabular}
\end{center}
\vspace{0.5cm}

\subsection{Implementação do Pipeline Completo}

\begin{lstlisting}[caption={Pipeline completo de matching},
                   label={lst:pipeline}]
class MotorDeMatching:
    """
    Pipeline integrado: Pontuacao Ponderada -> Cosine Similarity.
    """

    def __init__(self, vocabulario: list[str],
                 limiar_ponderado: float = 0.4):
        self.vocabulario = vocabulario
        self.limiar_ponderado = limiar_ponderado  # Corte na etapa 1

    def rankear(self,
                candidatos: list[dict],
                vaga: dict,
                top_n: int = 10) -> list[dict]:
        """
        Executa o pipeline completo e retorna os top-N candidatos.
        """
        # ETAPA 1: Filtragem por Pontuacao Ponderada
        candidatos_filtrados = []
        for c in candidatos:
            r = calcular_score_ponderado(
                Candidato(c["id"], c["id"],
                          c["competencias"],
                          c.get("projetos", 0),
                          c.get("experiencia", 0.0)),
                Vaga(vaga["id"], "", vaga["competencias_config"])
            )
            # Descarta candidatos com ausencias obrigatorias
            if r["ausencias_obrigatorias"] == 0:
                candidatos_filtrados.append(c)

        if not candidatos_filtrados:
            return []

        # ETAPA 2: Ranking por Cosine Similarity
        ranking = construir_matriz_scores(
            candidatos_filtrados, vaga, self.vocabulario
        )

        return ranking.head(top_n).to_dict(orient="records")

\end{lstlisting}



% --------------------------------------------------
\section*{Referências}

\begin{enumerate}
  \item SALTON, G.; MCGILL, M. J. \textit{Introduction to Modern Information Retrieval}.
        New York: McGraw-Hill, 1983.
  \item SCIKIT-LEARN DEVELOPERS. \textit{sklearn.metrics.pairwise.cosine\_similarity}.
        Disponível em: \url{https://scikit-learn.org/stable/modules/generated/sklearn.metrics.pairwise.cosine_similarity.html}
  \item PEDREGOSA, F. et al. Scikit-learn: Machine Learning in Python.
        \textit{Journal of Machine Learning Research}, v. 12, p. 2825–2830, 2011.
\end{enumerate}

\end{document}
